\subsection{Bounds on depth}

PROPOSITION 12.--- Let A be a Noetherian local ring, M a nonzero finitely generated A-module, and J an ideal of A distinct from A. We have the sequence of inequalities

\[
\begin{array}{r l} \operatorname{prof} _ {A} (J; M) & \leqslant \operatorname{codim} (\operatorname{Supp} (M) \cap V (J), \operatorname{Supp} (M)) \leqslant \dim (M) - \dim (M / J M) \\ & \leqslant [ J / \mathfrak {m} _ {A} J: \kappa_ {A} ]. \end{array}
\]

For every \(\mathfrak{p} \in \operatorname{Supp}(\mathrm{M})\cap \mathrm{V}(\mathrm{J})\), \(\operatorname{prof}_{\mathrm{A}}(\mathrm{J};\mathrm{M})\) is less than \(\dim_{\mathrm{A}_{\mathfrak{p}}}(\mathrm{M}_{\mathfrak{p}})\) (No. 5, Prop. 8 and No. 4, Cor. 2 of Th. 2), that is, (VIII, § 1, No. 4, Prop. 9)

to \(\operatorname{codim}(V(\mathfrak{p}),\operatorname{Supp}(\mathrm{M}))\). When \(\mathfrak{p}\) ranges over \(\operatorname{Supp}(\mathrm{M}) \cap \mathrm{V}(\mathrm{J})\), \(V(\mathfrak{p})\) describes the irreducible closed subsets of \(\operatorname{Supp}(\mathrm{M}) \cap \mathrm{V}(\mathrm{J})\), whence the first inequality. The second follows from Prop. 3 of VIII, § 1, No. 2. Moreover, one can find a generating set of J of cardinality \([J/\mathfrak{m}_{\mathrm{A}}J:\kappa_{\mathrm{A}}]\) (II, § 3, No. 2, Cor. 2 of Prop. 4); the third inequality then follows from VIII, § 3, No. 2, formula (8).

Remark 1.-- Consider the chain of inequalities in Prop. 12.

\begin{enumerate}
\def\labelenumi{\alph{enumi})}
\item
  For us to have \(\mathrm{prof}_{\mathrm{A}}(\mathrm{J};\mathrm{M}) = [\mathrm{J} / \mathfrak{m}_{\mathrm{A}}\mathrm{J}:\kappa_{\mathrm{A}}]\) it is necessary and sufficient that J can be generated by an M-regular sequence (no. 3, cor. 2 of th. 1 and A, X, p. 160, cor. 1).
\item
  The equality \(\dim(\mathrm{M}) - \dim(\mathrm{M}/\mathrm{JM}) = [\mathrm{J}/\mathfrak{m}_{\Lambda}\mathrm{J} : \kappa_{\Lambda}]\) means that J can be generated by a regular sequence for M (VIII, § 3, no. 2).
\item
  If M is Macaulay, we have \(\operatorname{prof}_{\mathrm{A}}(\mathbf{J};\mathbf{M}) = \dim (\mathbf{M}) - \dim (\mathbf{M} / \mathbf{JM})\) (§ 2, no. 2, cor. of prop. 2).
\end{enumerate}

Lemma 3.--- Let A be a Noetherian ring, \(\mathfrak{p} \subset \mathfrak{p}_1 \subset \ldots \subset \mathfrak{p}_{r-1} \subset \mathfrak{q}\) a saturated chain of length \(r\) of prime ideals of A, 'M a \(\Lambda\) -module of finite type and \(n\) an integer. If the A-module \(\mathrm{Ext}_{\mathrm{A}_{\mathfrak{p}}}^{n}(\kappa(\mathfrak{p}), \mathrm{M}_{\mathfrak{p}})\) is nonzero, the same holds for \(\mathrm{Ext}_{\mathrm{A}_{\mathfrak{q}}}^{n+r}(\kappa(\mathfrak{q}), \mathrm{M}_{\mathfrak{q}})\) .

It obviously suffices to treat the case \(r = 1\); substituting then A, M, p, and q by \(A_q\) , \(M_q\) , \(pA_q\) and \(qA_q\) respectively, we are reduced to treating the case where A is local and where \(q = m_A\) . Let \(x\) be an element of \(m_A - p\) . The \(A_p\) -module \(\text{Ext}_A^n(A/p, M) \otimes_A A_p\) is isomorphic to \(\text{Ext}_{A_p}^n(\kappa(p), M_p)\) (A, X, p. 111, prop. 10 b)), hence is nonzero by hypothesis; a fortiori \(\text{Ext}_A^n(A/p, M)\) is not zero. The exact sequence

\[
0 \rightarrow \mathrm{A} / \mathfrak {p} \stackrel {{x _ {\mathrm{A} / \mathfrak {p}}}} {{\longrightarrow}} \mathrm{A} / \mathfrak {p} \longrightarrow \mathrm{A} / (\mathfrak {p} + x \mathrm{A}) \rightarrow 0
\]

provides an exact sequence of extension modules

\[
\operatorname{Ext} _ {\mathrm{A}} ^ {n} (\mathrm{A} / \mathfrak {p}, \mathrm{M}) \xrightarrow {u} \operatorname{Ext} _ {\mathrm{A}} ^ {n} (\mathrm{A} / \mathfrak {p}, \mathrm{M}) \longrightarrow \operatorname{Ext} _ {\mathrm{A}} ^ {n + 1} (\mathrm{A} / (\mathfrak {p} + x \mathrm{A}), \mathrm{M}),
\]

where \(u\) is the homothety of ratio \(x\) . By Nakayama's lemma, this is not surjective, so the A-module \(\text{Ext}_A^{n+1}(\text{A}/(\mathfrak{p} + x\text{A}), \text{M})\) is not zero. But the only prime ideal of A containing \(\mathfrak{p} + x\text{A}\) is \(\mathfrak{m}_\text{A}\); therefore the A-module \(\text{A}/(\mathfrak{p} + x\text{A})\) has finite length (VIII, § 3, no. 2, lemma 2). If \(\text{Ext}_\text{A}^{n+1}(\kappa_\text{A}, \text{M})\) were zero, one would deduce, by induction on the length of N, the vanishing of \(\text{Ext}_\text{A}^{n+1}(\text{N}, \text{M})\) for every A-module N of finite length. This contradiction completes the proof.

PROPOSITION 13.--- Let A be a Noetherian ring, M a finitely generated A-module, and p and q two prime ideals in Supp(M) with p ⊂ q. Then

\[
\operatorname{prof} _ {A _ {\mathfrak {q}}} (M _ {\mathfrak {q}}) \leqslant \operatorname{prof} _ {A _ {\mathfrak {p}}} (M _ {\mathfrak {p}}) + \dim (A _ {\mathfrak {q}} / \mathfrak {p} A _ {\mathfrak {q}}).
\]

More precisely, for every saturated chain of prime ideals \(\mathfrak{p} \subset \mathfrak{p}_1 \subset \ldots \subset \mathfrak{p}_{r-1} \subset \mathfrak{q}\) , one has \(\operatorname{prof}_{\mathrm{A}_{\mathfrak{q}}}(\mathrm{M}_{\mathfrak{q}}) \leqslant \operatorname{prof}_{\mathrm{A}_{\mathfrak{p}}}(\mathrm{M}_{\mathfrak{p}}) + r\) .

Set \(p = \text{prof}_{\text{A}_p}(\text{M}_p)\) and let us prove the second inequality. It is evident if \(p = +\infty\) otherwise, we have \(\text{Ext}_{\text{A}_p}^p(\kappa(\mathfrak{p}), \text{M}_p) \neq 0\) , whence \(\operatorname{Ext}_{\Lambda_{\mathfrak{q}}}^{p+r}(\kappa(\mathfrak{q}),\mathrm{M}_{\mathfrak{q}})\neq0\) by lemma 3, which implies \(\operatorname{prof}_{\mathrm{A}_{\mathfrak{q}}}(M_{\mathfrak{q}})\leqslant p+r\) . Since \(\dim(A_{\mathfrak{q}}/\mathfrak{p}A_{\mathfrak{q}})\) is the supremum of the lengths of saturated chains of prime ideals with given endpoints. \(\mathfrak{p}\) and \(\mathfrak{q}\) , the first assertion is a consequence of the second.

COROLLARY 1.--- We have the inequality

\[
\dim (M _ {\mathfrak {q}}) - \operatorname{prof} _ {\Lambda_ {\mathfrak {q}}} (M _ {\mathfrak {q}}) \geqslant \dim (M _ {\mathfrak {p}}) - \operatorname{prof} _ {A _ {\mathfrak {p}}} (M _ {\mathfrak {p}}) \geqslant 0.
\]

This follows from prop. 13 and the inequality \(\dim(M_{\mathfrak{q}}) \geqslant \dim(M_{\mathfrak{p}}) + \dim(A_{\mathfrak{q}} / \mathfrak{p}A_{\mathfrak{q}})\) (VIII, § 1, no. 4, prop. 9, a) and no. 3, prop. 7, b)).

COROLLARY 2.--- Let A be a Noetherian local ring and M a finitely generated A-module. Then one has the inequality

\[
\operatorname{prof} _ {\mathrm{A}} (\mathrm{M}) \leqslant \inf _ {\mathfrak {p} \in \operatorname{Ass} (\mathrm{M})} \dim (\mathrm{A} / \mathfrak {p}).
\]

Let \(\mathfrak{p}\) a prime ideal associated with M; we have \(\mathrm{prof}_{\mathrm{A}_{\mathfrak{p}}}(\mathrm{M}_{\mathfrak{p}})=0\) (No. 1, Remark 2). Prop. 13 applied to the ideals \(\mathfrak{p}\subset\mathfrak{m}_{\mathrm{A}}\) provides the inequality \(\mathrm{prof}_{\mathrm{A}}(\mathrm{M})\leqslant\dim(\mathrm{A}/\mathfrak{p})\) , whence the corollary.

Remark 2.--- We have \(\sup_{\mathfrak{p}\in\operatorname{Ass}(\mathbb{M})}\dim(\mathrm{A}/\mathfrak{p})=\dim(\mathrm{M})\) (VIII, § 1, no. 4, remark 2) and one recovers the inequality \(\operatorname{prof}(\mathrm{M})\leqslant\dim(\mathrm{M})\) for \(\mathrm{M}\neq0\) (No. 4, cor. 2 of th. 2).

